% results_discussion

\section{Discussion}
\label{sec:discussion}

Across six research questions, one pattern recurs: satisfaction with a
devotional app is governed by the reliability of its core functions, not by
how many features it offers, and three framings borrowed from the broader HCI
and self-tracking literatures mis-locate where the problem sits. We take up
what each reframing changes, then what the paper implies methodologically for
research on app-store reviews beyond this domain. The prescriptive
consequences are collected separately in \S\ref{sec:design-implications}.

\subsection{Three Reframings}

\textbf{Gamification.} The harm the self-tracking literature predicts is guilt
induced by streaks and scores. That is not what we find (\S\ref{sec:rq1}). The
harm is a record users cannot trust and a prompt that does not arrive: an app
that loses logged history, regresses a count, or silently stops prompting
before logging is even at issue. Where a genuinely normative objection does
surface, it is not that the app makes the user feel judged; it is that the app
\emph{models worship incorrectly} --- no menstrual exemption, no support for
supererogatory (\emph{sunnah}) prayer, imposed goals that displace a user's own
targets. Softening a streak mechanic does not address a tracker that does not
know a user is exempt from prayer that day.

\textbf{Accuracy versus aesthetics.} The trade-off this framing poses is
mis-specified (\S\ref{sec:rq2}). Accuracy is not one thing to be weighed
against interface polish: prayer-time correctness behaves like a reliability
requirement, disjoint in cost from every interface complaint and carrying the
corpus's second-highest penalty, while madhab configuration behaves like an
ordinary preference, cheap to get wrong and reversing sign under a robustness
check. Treating ``accuracy'' as a single design dimension obscures exactly the
distinction that matters for prioritisation.

\textbf{Feature bloat.} What users call bloat is a property of what the
interface \emph{surfaces} in the moment, not of how many capabilities the
application \emph{contains} (\S\ref{sec:rq5}). The feature-count framing cannot
express that distinction, which is why it fails in both directions at once: no
relationship to bloat complaints where one is predicted, and an inverse one
where none is.

\subsection{The Two Gaps}

RQ6 (\S\ref{sec:rq6}) names two failures that the other five research questions
independently arrive at from different directions. Coverage is a \emph{market
failure}: whole categories of need are served by under a quarter of the
ecosystem despite demand that is explicit and easy to find in review text.
Delivery is a \emph{quality failure}, and the more consequential, because it
falls on promises applications have already made rather than on capabilities
they never claimed. The two most-broken promises in the corpus are adhan
reminders and monetisation terms, which is to say the core function of a
prayer-times app and the terms it is offered on.

\subsection{Methods Implications}

Two of this paper's four contributions are methodological, and they carry three
lessons for app-store review research beyond this domain.

\textbf{Star ratings are an unsafe proxy for defect severity.} 41.6\% of
reviews carrying a negative accuracy complaint award four or five stars
(\S\ref{sec:rq2}). A defect-prioritisation process that reads severity off the
star distribution --- the signal most app-store ranking and triage systems are
built around --- will systematically under-weight exactly the class of defect
users are least willing to let tank their rating. Text, not stars, is where
this class of complaint is visible.

\textbf{Single-app qualitative sampling can manufacture genre-level themes.}
Themes that looked like general properties of gamified prayer tracking
accounted for 51\% of a single-application coding pass and only 4\% of an
independently drawn, stratified cross-application pass on the same phenomenon
(\S\ref{sec:rq1-coding}). Both samples were coded in good faith; the difference
is entirely in what a selection rule's tie-break happened to return. Any
qualitative app-review study drawing its sample from one application, even
inadvertently, should treat its themes as a hypothesis about that application
until checked against others.

\textbf{Devotional register breaks off-the-shelf NLP.} A feature-vocabulary
match scored 0\% precision on its own gold rows because reviewers routinely
close devotional app reviews with formulae such as ``May Allah reward you,''
which generic goal-related keywords mistake for feedback about a tracking
feature (\S\ref{sec:method-devotional-filtering}). Anyone applying
general-purpose sentiment or aspect tooling to religious or otherwise
register-heavy text should expect this failure and budget for a filtering pass
before trusting aspect-level output.

\subsection{Three Threads}

Three concerns recur across research questions rather than belonging to any
one, and naming them is part of what makes six research questions a single
argument. \emph{Advertising and monetisation} run together through the results without
being the same thing: intrusive advertising is the costliest complaint in the
corpus by coefficient (\S\ref{sec:rq2}) and a dominant theme in bloat discourse
(\S\ref{sec:rq5}), while monetisation, tracked separately in our lexicon, is an
RQ1 coding theme and one of RQ6's two most-broken promises. \emph{Adhan reminders} are the most-cited reason users give for
choosing an app (\S\ref{sec:rq3}), one of the two most-broken promises
(\S\ref{sec:rq6}), and the dominant tracker-dissatisfaction theme at 46\%
(\S\ref{sec:rq1}). Beneath both sits a concern with \emph{the record rather
than the streak}: RQ1's reframing of gamification harm, RQ4's menstrual
exemption as a record the app gets wrong rather than a feature it lacks, and
RQ6's coverage gap are three instances of one underlying want --- an accurate,
owned, accommodating record of one's own worship.
